\phantomsection\addcontentsline{toc}{section}{Digital Twin}
\ECchapter{21}{Digital Twin}{Creating virtual replicas of physical systems}{ECBlue}

\ECObjectives{
\begin{itemize}
\item Explain how measurements continuously update a digital model.
\item Distinguish simulation, monitoring, diagnosis and prediction.
\item Identify the limits and uncertainties of a digital twin.
\end{itemize}}

\begin{multicols}{2}
\begin{engineeringnews}
Digital twins are increasingly used for aircraft engines, factories, bridges and energy systems. By
combining sensor data with physical or statistical models, engineers can estimate internal states,
predict degradation and test operating strategies without disturbing the real asset.
\end{engineeringnews}

\begin{ecarticle}
A digital twin is a virtual representation of a physical system that is regularly updated with real
measurements. Unlike a static CAD model, it evolves with the asset. Sensors measure quantities such as
temperature, vibration, force, pressure or electrical current. These data are filtered, time-stamped
and compared with a mathematical model that describes how the system should behave.

The model may be physics-based, data-driven or hybrid. A physics-based model uses conservation laws,
material properties and geometry. A data-driven model learns relationships from historical data. A
hybrid twin combines both approaches: physics constrains the solution while machine learning corrects
unknown effects or accelerates computation.

Engineers use twins for several linked purposes. Monitoring compares measured and expected behaviour.
A growing difference can reveal sensor drift, wear or an incorrect model. Diagnosis tries to identify
the cause of that difference. Prediction estimates future states, for example the remaining life of a
bearing or the battery temperature during a demanding mission. Optimisation allows engineers to test
control settings, maintenance dates or production schedules before applying them to the real system.

A useful twin depends on data quality. Measurements must be calibrated, synchronised and transmitted
reliably. The model must also be validated over the full operating range. A twin that performs well at
nominal speed may fail during start-up, overload or extreme weather. Its uncertainty must therefore be
reported rather than hidden.

Digital twins do not replace physical tests. They help engineers decide which tests are most useful and
interpret measurements more effectively. Their value comes from the continuous loop between the real
asset, the data platform and the model.
\end{ecarticle}

\begin{tcolorbox}[title=\sffamily\bfseries Did You Know?,colback=yellow!10,colframe=ECOrange]
A model becomes a true digital twin only when information flows repeatedly between the physical asset
and its virtual representation. Without this update loop, it remains a simulation model.
\end{tcolorbox}

\begin{engineeringfocus}
\textbf{Information loop}
\begin{center}
\resizebox{\linewidth}{!}{%
\begin{tikzpicture}[font=\scriptsize,>=Latex,node distance=4mm]
\node[draw=ECBlue,fill=ECLightBlue,rounded corners,align=center] (asset) {physical\\asset};
\node[draw=ECGreen,fill=ECGreen!10,rounded corners,align=center,right=of asset] (data) {sensors and\\data platform};
\node[draw=ECPurple,fill=ECPurple!10,rounded corners,align=center,right=of data] (twin) {digital\\twin};
\node[draw=ECOrange,fill=ECOrange!10,rounded corners,align=center,below=of data] (decision) {diagnosis and\\decision};
\draw[->,thick] (asset)-- node[above]{measurements} (data);
\draw[->,thick] (data)-- node[above]{updates} (twin);
\draw[->,thick] (twin)--(decision);
\draw[->,thick] (decision)--(asset);
\end{tikzpicture}%
}
\end{center}
\end{engineeringfocus}

\begin{keyfigures}
\begin{tabularx}{\linewidth}{@{}Xr@{}}
Essential input & calibrated data\\
Core operation & model update\\
Typical output & health estimate\\
Main risk & model drift\\
Validation method & test comparison
\end{tabularx}
\end{keyfigures}

\begin{tcolorbox}[title=\sffamily\bfseries Prediction error over time,colback=white,colframe=ECBlue]
\begin{center}
\begin{tikzpicture}
\begin{axis}[width=0.96\linewidth,height=46mm,xlabel={Time (days)},ylabel={Error (\%)},grid=major,
ticklabel style={font=\scriptsize}]
\addplot[thick,mark=*] table[x=time,y=error,col sep=comma]{data/EC21/prediction.csv};
\end{axis}
\end{tikzpicture}
\end{center}
\footnotesize A rising residual may indicate wear, a changing environment or an outdated model.
\end{tcolorbox}

\begin{applicationexercise}{Bearing health indicator}
A digital twin predicts a bearing vibration level of $2.4\,\mathrm{mm\,s^{-1}}$, while the sensor
measures $3.0\,\mathrm{mm\,s^{-1}}$.
\begin{enumerate}
\item Calculate the absolute residual.
\item Calculate the relative residual as a percentage of the predicted value.
\item The warning threshold is $20\%$. Decide whether maintenance should be investigated.
\end{enumerate}
\end{applicationexercise}

\begin{vocabulary}
\begin{tabularx}{\linewidth}{@{}lX@{}}
digital twin & jumeau numérique\\
asset & équipement réel\\
state estimation & estimation d'état\\
residual & résidu / écart\\
calibration & étalonnage\\
model drift & dérive du modèle\\
remaining useful life & durée de vie restante\\
uncertainty & incertitude\\
data-driven & fondé sur les données\\
validation & validation
\end{tabularx}
\end{vocabulary}

\begin{questions}
\item How is a digital twin different from a CAD model?
\item What is the role of sensors in the update loop?
\item Compare physics-based and data-driven models.
\item Why must a twin be validated outside nominal operation?
\item What can cause prediction error to increase?
\end{questions}

\begin{engineeringchallenge}
Design a digital twin for a school ventilation unit. Select four sensors, define the model outputs,
propose one fault indicator and explain how the twin could reduce energy use without lowering air
quality.
\end{engineeringchallenge}

\begin{speaking}
Should safety-critical decisions ever be made automatically by a digital twin? Discuss reliability,
responsibility, human supervision and uncertainty.
\end{speaking}
\end{multicols}

\subsection*{Sources}
\ECsource{NIST -- Digital Twins}{https://www.nist.gov/topics/digital-twins}
\ECsource{NASA -- Digital Twin}{https://ntrs.nasa.gov/}